\scanpage{69}
\mspage{2}
alors $\unsure{\mathbf{\Gamma}|\Gamma}$ opère non seulement sur \deleted{$X$} l'ens. $X$, mais aussi sur le groupoïde $\widehat\Pi_1(X)$, à cause du fait que celui-ci admet une définition purement schématique, ne faisant pas intervenir la topologie de $\mathbb C$. Bien entendu, si $\lambda:\deleted{x} x\to y$ dans $\Pi_1$, alors pour $u\in\Gamma$, $u(\lambda)$ n'est \unsure{p|pas} en général une flèche de $\Pi_1$ (de $u(x)$ vers $u(y)$) mais seulement de $\widehat\Pi_1$.

\deleted{Notation} Notation ${}^ux$ ($x\in X,u\in\unsure{\mathbf{\Gamma}|\Gamma}$) et \deleted{\unsure{``}}
\[
{}^u\lambda:{}^ux-{}^uy\qquad\text{si }\lambda:x\to y\quad(\text{dans }\widehat\Pi_1(X)).
\]
On a
\[
(6)\qquad\begin{cases}
{}^u(\operatorname{id}_x)=\operatorname{id}_{{}^ux}\qquad(x\in X)\\
{}^u(\mu\lambda)={}^u\mu\,{}^u\lambda
\end{cases}
\]
\diagram{\(\left(\begin{tikzcd}[ampersand replacement=\&,column sep=large]
x\arrow[r,"\lambda"]\&y\arrow[r,"\mu"]\&z\\
{}^ux\arrow[r,"{}^u\lambda"]\&{}^uy\arrow[r,"{}^u\mu"]\&{}^uz
\end{tikzcd}\right)\)}
Notre objet est d'essayer de ``déterminer'' autant que faire se peut cette action de $\unsure{\mathbf{\Gamma}|\Gamma}$. À vrai dire, si $\bar K$ est la clôture algébrique de $K$ dans $\mathbb C$, on s'intéressera surtout aux sous-groupoïdes de $\Pi_1(X)$ et $\widehat\Pi_1(X)$ induits sur l'ens d'objets $X_{\bar K}$, qu'on désignera par $\Pi_1(X_{\bar K}),\widehat\Pi_1(X_{\bar K})$ -- le deuxième est stable par l'action de $\unsure{\mathbf{\Gamma}|\Gamma}$, et si on pose
\[
(7)\qquad\Gamma=\operatorname{Aut}_{K\text{-cor}}(\bar K)=\operatorname{Gal}_{\bar K/K},
\]
de sorte que $\Gamma$ est un quotient de $\unsure{\mathbf{\Gamma}|\Gamma}$, l'action de $\unsure{\mathbf{\Gamma}|\Gamma}$ sur $\Pi_1(X_{\bar K})$ se fait à travers $\Gamma$. $\uparrow$
\margin{68}
